\chapter{A reference sheet and a teach-back check}
\label{chap:reference}

\section{The main biological words}

\begin{description}[style=nextline]
 \item[Allele] A version of the genetic locus being followed. Our model has
       two versions, A and B.
 \item[Haploid model] Each individual has one copy of the locus, so it has
       either A or B in this model.
 \item[Focal type] The type whose fate we follow. We call it A.
 \item[Resident type] The alternative type already present. We call it B.
 \item[Selection] A difference in expected reproductive success between
       the types. Here it enters through parent-selection weights.
 \item[Genetic drift] Random change in allele frequency due to finite
       random sampling. It occurs even when the weights are equal.
 \item[Neutrality] Equal reproductive weights. In our parameter convention,
       the advantage is zero.
 \item[Loss] No A individuals remain. The population itself still exists.
 \item[Fixation] Every individual in the modelled population is A.
 \item[Absorbing state] A state that cannot be left under the rules. With no
       mutation, zero A individuals and all A individuals are absorbing.
 \item[Well mixed] The rules do not assign spatial restrictions to parent
       or replacement selection.
 \item[Trajectory or path] The ordered sequence of states from one run.
 \item[Replicate] A new complete run from the specified initial condition,
       using further random draws.
\end{description}

\section{Read the mathematical notation in ordinary words}

In the table below, each symbol is followed by its meaning. The same letters
must retain the same definitions when you explain the project.

\begin{longtable}{p{0.21\textwidth}p{0.72\textwidth}}
\toprule
Symbol & Ordinary meaning \\
\midrule
\endhead
\(N\) & The total number of individuals. It stays constant during a run. \\
\(i\) & The current number of A individuals. It changes during a run. \\
\(i_0\) & The number of A individuals at the start. In the main fixation
            experiment, this is one. \\
\(N-i\) & The current number of B individuals. \\
\(x=i/N\) & The current fraction, or frequency, of A. \\
\(s\) & The relative reproductive advantage of A in this guide's convention.
          The main code accepts finite, nonnegative values. \\
\(r=1+s\) & The reproductive weight of one A individual when one B
              individual has weight one. \\
\(W_i=ri+N-i\) & The total reproductive weight of the population at count
                   \(i\). \\
\(T_i^+\) & The chance that the next event increases the A count by one. \\
\(T_i^-\) & The chance that the next event decreases the A count by one. \\
\(T_i^0\) & The chance that the next event leaves the A count unchanged. \\
\(u_i\) & The probability of eventual A fixation, starting from count
            \(i\). This concerns the whole future run. \\
\(d_i=u_i-u_{i-1}\) & The change in fixation probability between two
                       neighbouring starting counts. \\
\(\sigma=\log(1+s)\) & The log-fitness difference. It is useful for the
                        stable calculation and reference-paper comparison. \\
\(R\) & The number of repeated complete simulation runs. \\
\(K\) & The number of those runs in which A fixes. \\
\(\widehat p=K/R\) & The estimated fixation probability from the simulated
                      runs. The hat marks it as an estimate. \\
\(\mathbb{E}[Y]\) & The probability-weighted average of the random quantity
                      called \(Y\). \\
\(\Delta x\) & The change in allele frequency during one event. \\
\(\Pr(B\mid A)\) & The probability of a probability event called \(B\),
                    given that an event called \(A\) occurred. These letters
                    are event labels in this notation. \\
\(\sum\) & Add the terms specified by the expression's index and limits. \\
\(\approx\) & Approximately equal; an approximation or rounded value is
               being used. \\
\bottomrule
\end{longtable}

The symbols used inside a mathematical derivation and the variable names
used inside a function need not be identical. For example, the model's
function \code{wilson_interval(k, n)} uses \code{k} and \code{n} for a
success count and replicate count. In that function, \code{n} is not the
population size \(N\). The argument's meaning comes from the function's
definition and the value passed to it.

\section{Programming terms in plain language}

\begin{description}[style=nextline]
 \item[Variable] A name referring to a value. Assignment gives that name
       its current value.
 \item[Function] A named set of instructions that takes inputs and can
       return an output. Defining a function does not by itself run its body.
 \item[Argument] A value supplied when calling a function. Its position
       or keyword tells the function how to use it.
 \item[Return value] The result a function gives back to its caller.
 \item[Loop] An instruction to repeat a block. A \code{for} loop visits
       a specified sequence; a \code{while} loop continues while a condition
       remains true.
 \item[Conditional] An \code{if}, \code{elseif}, or \code{else} choice
       that selects instructions according to a condition.
 \item[Boolean] A truth value, either \code{true} or \code{false}.
 \item[Array or vector] An ordered collection of entries. Julia's first
       ordinary array index is one, not zero.
 \item[Index] A position used to select an entry from a collection.
 \item[Tuple] An ordered group of values. A named tuple also gives fields
       descriptive names, accessed with a dot.
 \item[Module] A container grouping related Julia definitions and controlling
       which names are exported for convenient use.
 \item[Standard library] A library supplied with Julia, such as
       \code{Random}, \code{Printf}, or \code{LinearAlgebra}.
 \item[Random-number generator] An algorithm that supplies the sequence of
       numbers used to sample random outcomes in the simulation.
 \item[Seed] A recorded initial setting for the random-number generator.
       Reproducibility also depends on the code and software versions.
 \item[Broadcasting] Applying an operation separately to the entries of a
       collection. A dot on an operator, such as \code{.^}, requests this.
 \item[String] Text held as a value, commonly written in double quotation
       marks. A filename or message can be a string.
 \item[CSV] A text-file format with comma-separated columns and one row per
       line in our controlled outputs.
 \item[Working directory] The folder the terminal treats as the current
       location for relative commands and paths.
 \item[Validation] Checking that supplied values meet a function's stated
       requirements before using them.
 \item[Test] A specified check that compares behaviour with an expected
       property or independently known answer.
\end{description}

\section{Connect mathematical ideas to the code}

\begin{longtable}{p{0.39\textwidth}p{0.54\textwidth}}
\toprule
Mathematical or modelling task & Julia function or script \\
\midrule
\endhead
Check whether the inputs describe an allowed population &
\code{validate_parameters} \\
Calculate the three next-event probabilities &
\code{transition_probabilities} \\
Turn a uniform draw into increase, decrease, or no change &
\code{next_count} \\
Store one history to a chosen display horizon &
\code{simulate_trajectory} \\
Run until loss or fixation and return its outcome &
\code{simulate_fixation} \\
Evaluate the exact fixation probability &
\code{fixation_probability} \\
Solve the backward equations independently &
\code{backward_probabilities} \\
Calculate an interval for a sampled success proportion &
\code{wilson_interval} \\
Generate the deterministic mean-field comparison &
\code{deterministic_trajectory} \\
Print the toy example and a small repeated experiment &
\code{scripts/demo.jl} \\
Generate the grid, analytical comparisons, paths, and metadata &
\code{scripts/run_study.jl} \\
Convert full-study outputs into document inputs &
\code{scripts/prepare_guide.jl} \\
Check specified implementation properties &
\code{test/runtests.jl} \\
\bottomrule
\end{longtable}

\section{Five distinctions to keep clear}

\begin{enumerate}
 \item \textbf{An event probability and a fixation probability.}
       The increase probability describes the next attempt. Fixation
       probability describes the eventual absorbing outcome.
 \item \textbf{A parameter and a state.} Population size and advantage
       are fixed settings. The current A count is the changing state.
 \item \textbf{A deterministic prediction and an average of simulations.}
       The deterministic curve follows a mean-field recurrence. It is not
       generally the exact average of the random trajectories.
 \item \textbf{An estimate and an exact answer.} The simulated proportion
       varies across finite samples. The exact probability is calculated
       from the model's mathematical rules.
 \item \textbf{A display horizon and an absorption stopping rule.}
       A displayed path can still be unfinished at its final plotted event.
       The complete fixation experiment continues until its outcome is known.
\end{enumerate}

\section{Teach the model back without looking}

Try answering the following questions aloud before reading the answer. You
have understood the project when you can connect each answer to a biological
rule, an equation, and a part of the code.

\begin{checkpoint}
\textbf{1. What does one model event do?}

\textbf{Answer.} It chooses a parent proportional to reproductive weight,
independently chooses a current individual uniformly for replacement, and
places a copy of the parent in that individual's place. The total number
of individuals stays fixed.

\textbf{2. Why multiply two fractions when calculating an increase?}

\textbf{Answer.} An increase needs both an A parent and a B replacement.
Given the current state, those selections are independent. Their joint
probability is the parent probability times the replacement probability.

\textbf{3. In the five-individual example, why is the chance of an A parent
one half even though only two individuals are A?}

\textbf{Answer.} With A weight \(1.5\), the total A weight is
\(2\times1.5=3\). The three B individuals have total weight three too.
A supplies three of the six weight units, so its parent probability is
\(3/6=1/2\).

\textbf{4. Why can the count stay unchanged after reproduction?}

\textbf{Answer.} An A offspring can replace A, or a B offspring can replace
B. These are real event attempts even though the type count is unchanged.

\textbf{5. Why can an advantageous allele be lost?}

\textbf{Answer.} A reproductive advantage makes an A parent more likely
relative to a B individual of the same abundance. It does not force A to
reproduce before being replaced. A sequence of chance events can reach zero
A individuals, especially when A begins rare.
\end{checkpoint}

\begin{checkpoint}
\textbf{6. How do you estimate fixation probability from the code?}

\textbf{Answer.} Start many complete runs at the same initial condition,
count how many end with all A, and divide that count by the number of runs.
Every replicate in this estimate runs until its absorbing outcome is known.

\textbf{7. What does a confidence interval add?}

\textbf{Answer.} It reports uncertainty from observing a finite number of
random outcomes. A nominal 95 percent interval procedure is designed for
approximately 95 percent coverage over repeated batches under its assumptions.
It does not assign a 95 percent probability to the fixed true parameter
after this particular interval has been calculated.

\textbf{8. Why is the neutral answer the initial frequency?}

\textbf{Answer.} In the backward-equation derivation, equal reproductive
weights make all neighbouring fixation-probability differences equal.
Starting from zero at count zero and ending at one at count \(N\), the
increment must be \(1/N\) per extra copy. At count \(i\), the answer is
therefore \(i/N\).

\textbf{9. What information lets us condition on just the next step?}

\textbf{Answer.} The current count and fixed parameters determine all future
transition rules. The previous path is not needed. This is the Markov
property. After the next event, we use the fixation probability associated
with the resulting count.

\textbf{10. Why is solving the backward equations a useful check?}

\textbf{Answer.} The solver obtains the probabilities from a set of linear
equations rather than evaluating the closed formula. Agreement tests whether
two different calculations of the same model match.
\end{checkpoint}

\begin{checkpoint}
\textbf{11. What does the deterministic curve leave out?}

\textbf{Answer.} It replaces random event outcomes with a smooth mean-field
change. It does not represent chance early loss from a positive initial
frequency. Its state can be fractional, whereas an actual A count is an
integer.

\textbf{12. Does twice the neutral fixation probability imply likely fixation?}

\textbf{Answer.} No. If neutral fixation probability is one percent, twice
that baseline is two percent. The relative improvement is substantial while
the absolute chance remains small.

\textbf{13. Why compare one initial copy with ten percent initial frequency?}

\textbf{Answer.} At one copy, increasing population size also lowers the
initial frequency. At a fixed frequency, the initial count grows with
population size. These are different questions and should not be mixed.

\textbf{14. Why should an approximation be checked against the exact result?}

\textbf{Answer.} An approximation simplifies an expression using assumptions
about its parameters. Those assumptions can fail. Exact-versus-approximate
error is a mathematical discrepancy, which more simulation replicates do
not remove.

\textbf{15. What would you say about the contribution of this project?}

\textbf{Answer.} It derives and reproduces established results for a stated
Moran model, validates a Julia implementation against an exact calculation,
and investigates selection, drift, and approximation accuracy over a defined
parameter range. Its purpose is to demonstrate understanding and careful
scientific computation. It is not a claim of a new biological discovery.
\end{checkpoint}

\section{A final practice exercise with a full answer}

\textbf{Task.} Explain the model to a friend in five sentences without using
the terms Markov chain, backward equation, or Monte Carlo.

\textbf{One possible answer.} We follow two genetic types in a population
whose size stays constant. At each step, one individual is chosen to
reproduce and its offspring replaces another randomly selected individual.
One type can have an advantage in the reproduction draw, but chance can
still eliminate it while it is rare. We repeat the process many times and
record how often that type eventually fills the population. We check the
computer results against a probability calculated directly from the same
rules.

Now give the technical version. Name the state, parameters, transition
probabilities, absorbing boundaries, fixation-probability recurrence, and
validation calculation. Being able to move between the ordinary and
technical explanations is a stronger sign of understanding than being able
to recite either version alone.
